\subsection{감사 AUD1: CMM 절}\label{AUD1:sec-CMM}

\paragraph{범위와 방법.}
감사 대상은 \texttt{agents/CMM/frag.tex}(합성 법 혼합 CMM 절)이며, 편집 전 원본은 \texttt{agents/CMM/frag.tex.orig\_before\_AUD1}에 보존하였다.
\PROVED\ 또는 \REFUTED\ 태그가 붙은 모든 항목(및 그 근거가 되는 \NUMERIC\ 항목)을 처음부터 다시 유도하였다:
각 부등식, 조건부 기댓값, 상수와 범위를 손으로 재계산하였고, CMM의 코드와 독립적으로 새로 작성한 코드로 수치를 재검증하였다.
\begin{itemize}[leftmargin=1.5em]
\item \texttt{agents/AUD1/code/aud1\_exact.py} $\to$ \texttt{out/aud1\_exact.txt}: $X\bmod q$의 정확한 법칙($U$는 부분집합 직접 열거, 또는 \emph{고정된} $U$),
 $\mu(b/q)$ 직접 합, 자체 구현한 $(\Z/d)^\times$ 지표 기계(2-멱 성분 $\langle-1\rangle\times\langle5\rangle$ 포함)로 명제 \texttt{CMM:prop-exact}의 상계,
 보조정리 \texttt{CMM:lem-good}, \texttt{CMM:lem-moment}의 점검, 켤레 검사.
\item \texttt{aud1\_constants.py} $\to$ \texttt{out/aud1\_constants.txt}: $N_1$ 하한, $e(k,v)$ 표, $R'$, $c_*$, $\beta_L$, $\theta_A$, 모멘트 장벽의 위치, 나쁜 이차 지표의 명시적 구성.
\item \texttt{aud1\_local.py} $\to$ \texttt{out/aud1\_local.txt}: 보조정리 \texttt{CMM:lem-local}을 정확한 $F_p$로 점검.
\item \texttt{aud1\_parseval.py} $\to$ \texttt{out/aud1\_parseval.txt}: 아래 명제~\ref{AUD1:prop-parseval}의 수치 확인.
\end{itemize}
CMM 파일의 모든 편집은 \texttt{\% AUD1:}로 시작하는 주석(이유 포함)으로 표시하였고, 편집 후 CMM 파일은 \texttt{test\_frag.sh}에서 \texttt{COMPILE OK}이다.

\paragraph{총평.}
판정 수: CONFIRMED 19, MINOR-FIX 6, GAP 0, WRONG 0.
핵심 수학(성분별 정확 전개, $2k$차 모멘트에 의한 나쁜 지표 계수, 곱셈적 상계, 점근식, 확대 singular series, 두 장벽 명제)은 모두 그대로 성립한다.
수정은 모두 범위, 해석, 표기에 관한 것이다. 가장 중요한 두 가지는 다음과 같다.
(i) 명제 \texttt{CMM:prop-badexist}(나쁜 원시 이차 지표의 존재)는 참이지만, CMM이 (O1)에서 끌어낸 ``따라서 $B(S)$들 사이의 상쇄가 필요하다''는 결론은
그 명제에서 나오지 않는다. 그런 지표 하나의 기여는 $\le\sqrt{q_S}/\varphi(q_S)$에 불과하다.
절댓값 전개가 \emph{실제로} 실패하는 지점은 새 Parseval 하한(명제~\ref{AUD1:prop-parseval}, \PROVED)이 준다:
$\log q\ge(\log2+o(1))L/\log L$이면 명제 \texttt{CMM:prop-exact}의 오차 상계는 $\ge2^{-\omega}$이다.
(ii) 주의 4.27(2)(d)의 ``격상''은 $q\le L^{\omega_0}\ll R_*$ 영역, 즉 무게가 이미 극소인 곳에만 해당하고, 실질적인 부분 $q\ge R_*$는 \OPEN\ 이다.
이 두 점과 그 밖의 범위 한정을 CMM 파일에 반영하였다.

%----------------------------------------------------------------------
\subsubsection{판정표}\label{AUD1:ssec-verdict}
\begin{center}\small
\begin{longtable}{p{0.25\textwidth}p{0.10\textwidth}p{0.13\textwidth}p{0.44\textwidth}}
\toprule
항목 (CMM 라벨) & 원 태그 & 판정 & 이유\\\midrule
\endhead
보조정리 \texttt{lem-N1} & PROVED (F2) & CONFIRMED & $N_1\ge\pi(L/2)-\pi(\sqrt L)-\omega$, Rosser--Schoenfeld 두 부등식, $\sqrt L\ge21.06$. 수치: $444\le L\le2\cdot10^5$에서 실패 없음(최소 비 $1.48$); 수치로는 $82\le L\le2\cdot10^5$ 전체에서 성립.\\
보조정리 \texttt{lem-fourier} & PROVED (F1) & CONFIRMED & $\hat f(\psi)=\psi(b)\tau(\overline\psi)/\varphi(p^m)$; 주지표는 Ramanujan 합 $c_{p^m}(1)=-1$ ($m=1$), $0$ ($m\ge2$); 비원시 비주지표는 $w=w_0+p^{m-1}h$ 분해로 $\tau=0$; $p=2$, $m=1$은 원시 지표가 없고 $c_1=-1=e(b/2)$로 일관.\\
주의 \texttt{rem-conj} & NUMERIC (켤레 없는 식 반박) & CONFIRMED & 켤레 판은 직교성으로 즉시 증명됨. 독립 수치: $q'=15,77,45,28$에서 켤레 판 오차 $\le8.5\cdot10^{-14}$, 켤레 없는 판 오차 $1.90$--$2.00$.\\
명제 \texttt{prop-exact} & PROVED & CONFIRMED & 부분분수, $J=\{e_j\ge k_j\}$로 조건화, $w_j=X/p_j^{e_j}$는 $p_j$와 서로소($c'$는 $p_j$를 안 가짐, $U$의 소수는 $>L/2\ge p_j$이므로 $q_j\notin E$), CRT로 원시 지표 분해, $|\hat G|=\sqrt d/\varphi(d)$, $c'$, $U$, $\mathbf e$의 독립성, $|\mathbb E\chi(U)|\le1$(따라서 $E$의 법칙 임의). 주항과 가중치 $\prod_{S}\frac1{v_j+1}\prod_{S^c}\rho^0_j$ 재계산 일치. 수치: 19개 사례(소수 멱, 2-멱, 고정 $U$ 포함)에서 참 오차 $\le$ 상계.\\
따름정리 \texttt{cor-mult} & PROVED & CONFIRMED & 곱의 전개와 $|M(q)|\le\prod\rho^0_j$.\\
보조정리 \texttt{lem-good} & PROVED & CONFIRMED & $|\tfrac{1+z}2|^2=\tfrac{1+\mathrm{Re}z}2\le\e^{-(1-\mathrm{Re}z)/2}$; $P'$의 소수는 $v_p=1$. 수치 위반 0.\\
보조정리 \texttt{lem-moment} & PROVED & CONFIRMED & 모든 $\chi\bmod d$에 대한 직교성, 잉여류 안의 $n\le(L/2)^k$ 개수 $\le(L/2)^k/d+1$, 순서배열 $\le k!$. 수치($k=1..4$) 위반 0.\\
명제 \texttt{prop-B} & PROVED & CONFIRMED & $k=\lceil\log d/\log(L/2)\rceil$이면 $(L/2)^k\ge d$이므로 $B_{\rm bad}\le2k!\kappa^k/\sqrt d$.\\
정리 \texttt{thm-main} (곱셈적 상계) & PROVED & CONFIRMED & 좋은 부분 $y_j^{1/2}\e^{-\eta N_1/(4r)}\le\eta_L$ ($N_1/(4\omega)\ge\sqrt L/12$); $x_j\le\frac32$, $k\le\min(x+1,2r)$, $(2r\kappa_*)^{x_j}=y_j^{\lambda_r}$; 경우 (a) $4r\kappa_*\le\e^r$, (b) $4r\kappa_*<Z^\mu<d^\mu$. 작은 소수에서 인수가 $1$을 넘을 수 있음은 CMM 스스로 명시(주의 \texttt{rem-size}); 해석은 주의~\ref{AUD1:rem-smallp}.\\
정리 \texttt{thm-main} ((H$_Z$-ii)) & PROVED (F3) & CONFIRMED & $\sum_b\chi(b)=0$로 주항 소거, $d\le Z\le(\log(L/2))^{A+1}$에서 SW, 둘째 합 $\le\pi(\sqrt L)+\omega$. 비유효.\\
정리 \texttt{thm-main} ($R_L$, $\theta_A\to1/A$) & PROVED & CONFIRMED & $\log t\le t/2+1$; 수치 $R_L=7$--$9\le13$--$15$, $\theta_A=0.458,0.388,0.327,0.288,0.238$ 재현.\\
따름정리 \texttt{cor-large} & PROVED (유효) & CONFIRMED & $c_r=(2r!)^{1/r}$ 비감소($(r+1)^r\ge2r!$), $\sqrt{q_S}>L^{r/4}$, $\lambda'\le\frac12$, $p_j>(L/2)^{1/2}$. F2만 사용(SW 불필요). $\beta_L\to0\iff\omega\log L=o(\sqrt L)$ 옳음. 수치 $\beta_L(10^{16})=0.0255$, $\beta_L(10^{12})=0.191$ 재현.\\
정리 \texttt{thm-asym} & PROVED (유효) & CONFIRMED & $\rho^0_j\le\frac12\e^{1/(y-1)}$, $\sum_{S\ne\emptyset}\sqrt{q_S}\le\prod(1+\sqrt{p_j})\le L^\omega$, $\binom\omega rr!\le\omega^r$, 비 $\le\frac12$.\\
따름정리 \texttt{cor-T3} & PROVED & MINOR-FIX & 부등식은 옳음($y=L/3$). ``$h'\le c\sqrt L/\log L\Rightarrow\ge2^{-h'}(1-13c-o(1))$''에서 가정이 $c\le0.156$을 강제하고 비자명성은 $c<0.078$임을 추가.\\
보조정리 \texttt{lem-local} & PROVED & CONFIRMED & $0.9\sigma\log p\le s_0/(v+1)<\sigma\log p$, $e(k,v)$ 표 재계산(최댓값 $-0.08985$, $(1,8)$), $(1,1)$ 지수 $\le-\delta$, 36쌍. 수치: (i) 비 $\le0.73$, (ii) 비 $\le4.8\cdot10^{-3}$.\\
정리 \texttt{thm-SS}(a) & PROVED (유효) & MINOR-FIX & 증명 옳음($\rho^{s_0}\le L^{-2-\delta/2}$, $\sum_{p\le L/2}p\le0.3138L^2/\log(L/2)$). ``$L_1$ 유효''는 $o(\cdot)$의 속도가 명시적일 때만이라는 단서 추가.\\
정리 \texttt{thm-SS}(b) & PROVED (유효) & CONFIRMED & $\Gamma_p\le\Gamma^\flat(p)$, 조건 (1)--(4) 모두 명시적; $p>L^{1/9}$는 보조정리 (ii), $p_*\le p\le L^{1/9}$는 $F_p\le1.5p^{-1.3}$. SW 불필요.\\
정리 \texttt{thm-SS}(c) & PROVED (F3) & CONFIRMED & $p<p_*$에서 $F_p\le1.5p^{4.5\Gamma^\flat(2)}$ (유계 인수, $v_p\to\infty$ 덕분), $1.5(\zeta(1.3)-1)=4.398$, 꼬리 $\le(\sum_{p>z}F_p)\prod(1+F_p)$.\\
주의 \texttt{rem-threshold} & PROVED & MINOR-FIX & $h$층 하한은 옳음. ``문턱은 합성 법을 포함해도 불변''을 정리 SS의 족 ($\omega(q)\le\omega_0$) 안으로 한정; 모든 $q\mid Q$는 \OPEN.\\
따름정리 \texttt{cor-upgrade}(1) & PROVED (F3) & CONFIRMED & 정리 SS(b)(c)의 재진술.\\
따름정리 \texttt{cor-upgrade}(2) & PROVED & MINOR-FIX & 부등식은 옳으나 정리 SS(c) 경유이므로 (F3) 표기 누락; 또 $q\le L^{\omega_0}\ll R_*$ 영역만 다루므로 4.27(2)(d)의 실질적 부분은 (3) \OPEN\ 임을 명시(요약·원장도 수정).\\
명제 \texttt{prop-barrier} & PROVED (방법 진술) & CONFIRMED & $\Phi_k=k!\kappa^k(d^{-1/2}+d^{1/2}(L/2)^{-k})$, $k<k_0$에서 비 $\ge(\frac{L/2}{k_0\kappa})^{k_0-k}\ge1$, $k_0!\ge(k_0/\e)^{k_0}$. 수치: $\min_k\Phi_k\ge1$이 되는 첫 $x$는 $(2.08$--$2.14)\sqrt L/\log L$ ($L=10^6,10^8,10^{10}$).\\
명제 \texttt{prop-badexist} & PROVED & CONFIRMED & $\mathbb F_2$ 선형대수, 법 소수의 비자명 지표는 원시. 수치 예: $L=1000$, $\omega=43>N_1=41$에서 $|S|=20$인 나쁜 이차 지표 구성.\\
요약·(O1)의 badexist 해석, 원장 & -- & MINOR-FIX & ``나쁜 지표가 있으므로 $B(S)$들 사이의 상쇄가 필요''는 따라나오지 않음(위 예에서 $\mathbb E\chi(c')=0$ 정확히). 올바른 장애는 명제~\ref{AUD1:prop-parseval}; 해당 문장을 \HEUR\ 로 바꾸고 이유를 교체.\\
\S\texttt{ssec-num} 표 & NUMERIC & MINOR-FIX & 모든 행을 독립 재현. $L=60$, $11\cdot13\cdot17$의 \#bad $225\to224$ ($\mathrm{Re}A=N_1/2$가 정확히 되는 경계 지표 2개의 부동소수점 문제); $8\cdot9\cdot5$ 행의 비어 있던 상계를 AUD1 값으로 채움.\\
\bottomrule
\end{longtable}
\end{center}

%----------------------------------------------------------------------
\subsubsection{lead가 지정한 중점 사항의 점검}\label{AUD1:ssec-focus}

\paragraph{(1) 지표/Gauss 합 전개.}
켤레는 계수 $\hat f(\psi)=\psi(b)\tau(\overline\psi)/\varphi(p^m)$에 정확히 들어가 있으며, 증명은 $|\hat G(\chi)|$만 쓰므로 켤레 누락 문제(p1p4 주의 10.4)가 생길 수 없다.
비원시 지표는 각 성분에서 $\tau=0$이 되어 사라지고, 주지표 성분은 Ramanujan 합 $-1$ ($m=1$) 또는 $0$ ($m\ge2$)으로 주항 $M(q)$의 인수
$\frac{v+1-k-1/(p-1)}{v+1}$을 만든다(소수 멱에서 $e_j=k_j-1$일 때만 $-1/(p-1)$이 기여). lead 스케치의 $J=\{j:q_j\mid c\}$로의 조건화는
CMM의 $J=\{e_j\ge k_j\}$와 같고, 조건화는 $c'$, $U$의 법칙을 바꾸지 않는다(지수의 독립성). $U(E)$의 법칙이 임의여도 되는 이유는 $|\mathbb E\chi(U)|\le1$로 버리기 때문이며,
고정된(결정론적) $U$로도 수치 확인하였다(\texttt{out/aud1\_exact.txt}의 ``U FIXED'' 행: 참 오차/상계 $=0.25$--$0.32$). $q_j\notin E$는 $q_j\le L/2<\min P$에서 옳다.
흥미로운 부수 관찰: $L=60$에서 $q=45,28$은 상계가 정확히 $0$이고 실제로 $\mu(a/q)=M(q)$(부동소수점 오차 내)이다. 예컨대 법 $9$에서는 $2$가 원시근이고 $v_2+1=6$이므로
모든 원시 $\chi\bmod9$에 대해 $\mathbb E\chi(2)^{e_2}=\frac16\sum_{e=0}^5\chi(2)^e=0$이다. 이는 정확 전개의 좋은 독립 검증이다.

\paragraph{(2) $2k$차 모멘트, $k$의 선택, 나쁜 지표 문턱.}
$\sum_{\chi\bmod d}|A(\chi)|^{2k}=\varphi(d)\mathcal N_k$는 $P'$의 소수가 $d$와 서로소이므로 성립하고, $\mathcal N_k\le N_1^kk!((L/2)^k/d+1)$은 곱의 크기 $\le(L/2)^k$와 소인수분해의 유일성에서 나온다.
$k=\lceil\log d/\log(L/2)\rceil$은 $(L/2)^k/d+1\le2(L/2)^k/d$를 보장하는 최소 선택이다. 문턱 $\mathrm{Re}A>(1-\eta)N_1$과 좋은 지표의 상계 $\e^{-\eta N_1/4}$은 일관된다
($\eta=\frac12$이 lead의 $|A|\ge N_1/2$, $\e^{-N_1/8}$). 19개 사례 모두에서 보조정리 \texttt{lem-good}와 모멘트 상계($k=1,\dots,4$)의 위반은 없었다.

\paragraph{(3) 상수와 범위, 유효성.}
$\kappa\le\kappa_*=\frac{3\log L}{2(1-\eta)^2}$ ($N_1\ge L/(3\log L)$), $\eta_L=\sqrt L\e^{-\eta\sqrt L/12}$ ($N_1/(4\omega)\ge\sqrt L/12$), $\eta'_L=\exp(\sqrt L/3-\eta L/(12\log L))$
($\omega\log L\le\sqrt L/3$), $\beta_L(\omega)$ (따름정리 \texttt{cor-large})의 정의와 쓰임이 모두 일치한다. 범위 $\omega=o(\sqrt L/\log L)$은 정확히 $\e(\omega\kappa_*)^{1/2}(L/2)^{-1/4}\to0$의 조건이고,
점근식의 범위 $\omega\kappa_*\le\frac12\sqrt y$ ($y=\sqrt L$이면 $\omega\le L^{1/4}/(2\kappa_*)$)도 옳다. ``유효'' 주장(따름정리 \texttt{cor-large}, 정리 \texttt{thm-asym}, 따름정리 \texttt{cor-T3}, 정리 SS(a)(b))은
F1(Gauss 합)과 F2(Rosser--Schoenfeld)만 쓰며 Siegel--Walfisz를 쓰지 않는다: 모든 $d$가 $>(\log L)^A$ 또는 $>\sqrt L$인 소인수를 가져 모멘트 계수만으로 충분하기 때문이다.
F3는 정리 \texttt{thm-main}의 (H$_Z$-ii), 정리 SS(c), 따름정리 \texttt{cor-upgrade}(1)(2)에서만 쓰인다((2)의 표기를 보충함).

\paragraph{(4) singular series 정리.}
$a$에 대한 합은 $\varphi(q)\le q$개의 항을 주고, $q\mid Q$ 위의 합은 곱셈적이어서 $\prod_p(1+F_p)$로 묶인다. 문턱 $s\ge(2+\delta)\log_2L$은 $(k,v)=(1,1)$ 지수
$2-\sigma\log2\le-\delta$에서 나오고, 나머지 35쌍은 $\sigma_0=2/\log2$에서 $\le-0.08985$이다(재계산). 꼬리 $\sum_{p>z}F_p\prod(1+F_p)$와 $\zeta(1.3)$ 상수도 옳다.
작은 $p$에서 $\rho_{p^k}>1$이 가능하지만 $v_p\to\infty$ 때문에 $\rho_{p^k}^{s_0}\le p^{\sigma(\Gamma_p-k)}$로 $L$과 무관하게 유계이므로 (c)의 $O_\eps(1)$은 정당하다.

\paragraph{(5) 장벽 명제.}
명제 \texttt{prop-barrier}의 부등식 사슬은 옳고(가정 $k_0\kappa\le L/2$ 필요), 수치로 장벽의 위치가 $\approx2.1\sqrt L/\log L$임을 확인했다.
명제 \texttt{prop-badexist}도 옳다. 그러나 그 의미는 ``모든 원시 지표가 좋다'' 경로의 불가능성에 국한되며, 모멘트 방법이나 절댓값 전개의 실패를 뜻하지 않는다:
$L=1000$, $\omega=43$에서 구성한 나쁜 이차 지표는 $|S|=20$, $\sqrt{q_S}/\varphi(q_S)\approx10^{-21}$이고, 작은 소수 $7$ ($\chi(7)=-1$, $v_7=3$) 때문에 $\mathbb E\chi(c')=0$이다.
절댓값 전개의 진짜 한계는 다음 명제가 준다.

\begin{proposition}[절댓값 전개의 Parseval 하한]\label{AUD1:prop-parseval}
\PROVED\ $q=p_1\cdots p_\omega$ 제곱 인수 없음, $p_j\in(\sqrt L,L/2]$ (따라서 $v_j=1$), $(a,q)=1$이라 하고, 명제 \texttt{CMM:prop-exact}의 오차 상계를
$\mathfrak B(q):=\sum_{\emptyset\ne S}B(S)\,2^{-|S|}\prod_{j\notin S}\rho^0_j$로 적자. $Q':=\prod_{p\mid Q,\,p\nmid q}p^{v_p}$, $\tau$는 약수의 개수이다. 그러면
\[
\mathfrak B(q)\ \ge\ 2^{-\omega}\,q^{-1/2}\Bigl(\frac{\varphi(q)}{\tau(Q')}-1\Bigr).
\]
특히 $\varphi(q)\ge2\tau(Q)\sqrt q$이면 $\mathfrak B(q)\ge2^{-\omega}\ge M(q)$, 즉 명제 \texttt{CMM:prop-exact}는 곱셈적 상계는 물론 $\mu(a/q)\approx M(q)$조차 줄 수 없다.
$\log\tau(Q)=(\frac{\log2}2+o(1))\frac L{\log L}$이고 $\varphi(q)/q\ge\frac13$ ($L$이 클 때)이므로 이 조건은 $\log q\ge(\log2+o(1))L/\log L$에서 성립한다
(뒤의 두 점근은 \EXT: PNT와 Mertens 곱 정리, Montgomery--Vaughan 2장과 6장; 교과서 사실, recalled).
% AUD1: theorem numbers deliberately not cited (not verified from source).
\end{proposition}
\begin{proof}
$\rho^0_j=\frac12(1+\frac1{p_j-1})\ge\frac12$이므로 모든 가중치가 $\ge2^{-\omega}$이다. $|\mathbb E\chi(c')|\le1$이므로
$B(S)=\frac{\sqrt{q_S}}{\varphi(q_S)}\sum_{\chi\ \text{원시}\bmod q_S}|\mathbb E\chi(c')|\ge q_S^{-1/2}\sum_{\chi\ \text{원시}\bmod q_S}|\mathbb E\chi(c')|^2\ge q^{-1/2}\sum_{\chi\ \text{원시}\bmod q_S}|\mathbb E\chi(c')|^2$.
$q$가 제곱 인수가 없으므로 법 $q$의 각 지표 $\chi$는 유일한 $S$에 대해 법 $q_S$의 원시 지표 $\chi^*$에서 유도되고, $(c',q)=1$이라 $\chi(c')=\chi^*(c')$이다.
$S=\emptyset$은 주지표에 해당하고 $\mathbb E\chi_0(c')=1$이다. 따라서
$\sum_{S\ne\emptyset}\sum_{\chi^*\ \text{원시}\bmod q_S}|\mathbb E\chi^*(c')|^2=\sum_{\chi\bmod q}|\mathbb E\chi(c')|^2-1=\varphi(q)\sum_r\mathbb P(c'\equiv r)^2-1\ge\frac{\varphi(q)}{N}-1$,
여기서 $N$은 $c'\bmod q$의 지지집합의 크기이고 마지막은 Cauchy--Schwarz이다. $c'$은 $Q'$의 약수이므로 $N\le\tau(Q')$. 합치면 주장이 나온다.
특히 부분: $\tau(Q')\le\tau(Q)$이고 $\varphi(q)\ge2\tau(Q)\sqrt q$이면 괄호 $\ge2\sqrt q-1\ge\sqrt q$이므로 $\mathfrak B(q)\ge2^{-\omega}$; $M(q)=\prod(\frac12-\frac1{2(p_j-1)})\le2^{-\omega}$.
점근: $\log\tau(Q)=\sum_{p\le L/2}\log(v_p+1)\le\pi(L/2)\log2+\pi(\sqrt L)\log(1+\log_2L)$와 PNT; $\varphi(q)/q\ge\prod_{\sqrt L<p\le L/2}(1-\frac1p)\to\frac12$ (Mertens).
\end{proof}
수치 확인 \NUMERIC\ (\texttt{code/aud1\_parseval.py}, \texttt{out/aud1\_parseval.txt}): $L=40$, $q=7\cdot11\cdot13\cdot17\cdot19$ ($N=\tau(Q')=72$)에서
$2^\omega\mathfrak B=80.7\ge2^\omega\cdot$(하한)$=5.06>1$; 다른 두 사례에서도 부등식 성립.
결론: 소수 $\in(\sqrt L,L/2]$로 된 $q$에 대해 명제 \texttt{CMM:prop-exact}의 절댓값 전개는 $\omega\ge(\log2+o(1))L/(\log L)^2$에서 원리적으로 쓸모없다.
그 아래 $\sqrt L/\log L\lesssim\omega\lesssim L/(\log L)^2$에서 절댓값 전개가 충분한지는 \OPEN\ 이다(모멘트 방법의 장벽은 $x=\log q_S/\log(L/2)\approx2.1\sqrt L/\log L$).
무게가 $1$인 $q\ge R_*$ 중 소인수가 모두 $(\sqrt L,L/2]$에 있는 것은, $\log R_*\ge c_RL/\log L$에서 $c_R>\log2$이면 전부 이 Parseval 영역에 들어간다(일반 $q\mid Q$로의 확장은 하지 않았다).

\begin{remark}[작은 소인수에 대한 곱셈적 상계의 해석]\label{AUD1:rem-smallp}
\PROVED\ 정리 \texttt{CMM:thm-main}에서 $\lambda,\theta\le\frac12$이고 $\eta_L\le1$이면 $\beta(y)\le\e+2$이므로, 제곱 인수가 없는 성분에서
$\rho_p-\pi_p=\frac{1/(p-1)+\beta(p)}{v_p+1}\le\frac{\e+3}{v_p+1}<\frac{6\log p}{\log L}$이다. 따라서 작은 $p$에서 인수 $\rho_p$가 $1$을 넘을 수 있더라도(CMM 주의 \texttt{rem-size})
$|\mu|\le\prod_p(\pi_p+\eps_p)$ 형태에서 $\eps_p=O(\log p/\log L)$이며, $p\ge L^{c}$에서는 $\eps_p\le\frac12(\frac1{p-1}+\e L^{-c\eps}+L^{c(\theta-1/2)}+\eta_L)$이다
($\lambda\le\frac12-\eps$일 때). 즉 CMM 요약의 ``$\omega(q)\le L^{1/2-\eps}$에서 CMM 증명''은 p1p4 주의 4.21의 형태 $\prod(\pi_p+\eps)$로 읽어도 정당하다(단 $\eps$는 작은 $p$에서 $\log p/\log L$ 크기).
\end{remark}

%----------------------------------------------------------------------
\subsubsection{CMM 파일에 가한 수정}\label{AUD1:ssec-edits}
\begin{enumerate}[label=(E\arabic*),leftmargin=2.5em]
\item 요약 목록: 주의 4.27(2)(d)의 격상 범위($q\le L^{\omega_0}\ll R_*$만)를 명시; 명제 \texttt{prop-badexist}가 방법의 장애가 아님을 한 문장으로 추가.
\item 따름정리 \texttt{cor-T3}: $c\le0.156$ (가정에서 강제), 비자명성 $c<0.078$을 추가. 원장 행도 수정.
\item 정리 \texttt{thm-SS}(a): ``$L_1$ 유효''에 ``$o(\cdot)$의 속도가 명시적이면'' 단서 추가.
\item 주의 \texttt{rem-threshold}: ``문턱 불변''을 $\mathcal Q_1$--$\mathcal Q_3$ ($\omega(q)\le\omega_0$) 안으로 한정, 전체는 \OPEN.
\item 따름정리 \texttt{cor-upgrade}(2): ``(F3 사용, 정리 SS(c) 경유)'' 추가, 무게가 극소인 영역만 다룬다는 주의 추가. 원장 행을 ``$\omega(q)\le\omega_0$ 부분, (F3); $q\ge R_*$는 (3) \OPEN''으로 수정.
\item (O1): ``따라서 $B(S)$들 사이의 상쇄가 필요''를 삭제하고, 나쁜 지표 하나의 기여가 무시할 만함(수치 예 포함), 명제~\ref{AUD1:prop-parseval}의 Parseval 하한을 올바른 장애로 인용, 전역 양 $Q\,\mathbb P(Y\equiv R\bmod Q)$의 필요성은 \HEUR, 중간 구간은 \OPEN\ 으로 표기.
\item 수치 표: \#bad $225\to224$; $8\cdot9\cdot5$ 행 ($L=60,100$)에 AUD1이 계산한 명제 상계 $1.25\cdot10^{-1}$, $1.92\cdot10^{-3}$과 \#bad $=1$을 기입.
\end{enumerate}
태그의 격하(\PROVED$\to$\COND{}/\HEUR/\OPEN)는 필요하지 않았다: 모든 \PROVED\ 항목의 증명이 적힌 그대로 완결되어 있다.

%----------------------------------------------------------------------
\subsubsection{독립 수치 요약}\label{AUD1:ssec-num}
\NUMERIC\ (\texttt{agents/AUD1/code/aud1\_exact.py} $\to$ \texttt{agents/AUD1/out/aud1\_exact.txt}). 표의 값은 $2^\omega$배이다.
\begin{center}\small
\begin{tabular}{rllrrrr}
\toprule
$L$ & $q$ & $(|P_i|,t_i)$ & $2^\omega M$ & $2^\omega\max|\mu-M|$ & $2^\omega\mathfrak B$ & 참/상계\\
\midrule
60 & $29\cdot23$ & (4,2) & 0.9205 & $1.25\cdot10^{-2}$ & $9.30\cdot10^{-2}$ & 0.135\\
60 & $29\cdot23$ & $U$ 고정 & 0.9205 & $2.42\cdot10^{-2}$ & $9.30\cdot10^{-2}$ & 0.260\\
60 & $11\cdot13\cdot17$ & (3,1) & 0.7734 & $4.34\cdot10^{-2}$ & $4.39\cdot10^{-1}$ & 0.099\\
100 & $47\cdot43$ & (4,2) & 0.9550 & $1.29\cdot10^{-3}$ & $1.29\cdot10^{-2}$ & 0.100\\
100 & $3\cdot7\cdot11$ & (3,2) & 1.5400 & $8.11\cdot10^{-4}$ & $6.18\cdot10^{-3}$ & 0.131\\
60 & $8\cdot9\cdot5$ & (3,1) & 0.5833 & $2.18\cdot10^{-2}$ & $1.25\cdot10^{-1}$ & 0.174\\
100 & $8\cdot9\cdot5$ & (3,1) & 1.0000 & $4.63\cdot10^{-4}$ & $1.92\cdot10^{-3}$ & 0.241\\
60 & $16\cdot3$ & (3,2) & 0.4167 & $1.17\cdot10^{-2}$ & $4.86\cdot10^{-2}$ & 0.241\\
100 & $49\cdot11$ & (4,2) & 0.5000 & $5.59\cdot10^{-4}$ & $4.24\cdot10^{-3}$ & 0.132\\
100 & $27\cdot25$ & $U$ 고정 & 0.3000 & $1.66\cdot10^{-3}$ & $5.24\cdot10^{-3}$ & 0.318\\
300 & $149\cdot139$ & (4,2) & 0.9860 & $8.14\cdot10^{-8}$ & $5.69\cdot10^{-7}$ & 0.143\\
300 & $19\cdot23\cdot29$ & (3,1) & 0.8693 & $2.89\cdot10^{-7}$ & $1.05\cdot10^{-6}$ & 0.276\\
40 & $7\cdot11\cdot13\cdot17$ & $U$ 고정 & 0.6445 & $2.52$ & $9.38$ & 0.269\\
\bottomrule
\end{tabular}
\end{center}
관찰: CMM 표의 모든 공통 행이 유효숫자 범위에서 재현되었고, 소수 멱($2^4$, $7^2$, $3^3\cdot5^2$ 포함)과 고정 $U$에서도 주항 공식과 상계가 성립한다.
마지막 행($L=40$)은 작은 $L$에서 실제 상대오차가 $1$을 넘을 수 있음을 보여 주며(상계는 여전히 성립), 증명된 상수가 $L\gtrsim10^{12}$에서야 의미를 갖는다는 CMM의 정직한 서술과 일치한다.
상수 점검(\texttt{out/aud1\_constants.txt}): $e(k,v)$ 표, $R'=7$, $c_*=3.732$, $\kappa_*=1.852\log L$ ($\eta=0.1$), $\sigma\le4.492$, $2^{-1.596}=0.3308$, $4\sqrt3\cdot1.852=12.83$ 모두 CMM과 일치.
